\subsection{On-Policy Distillation for LLM Post-Training}

On-policy distillation evolves beyond SFT and RLVR by combining trajectory exploration with dense, token-level guidance~\cite{lu2025onpolicydistillation,fu2025pruningweightstruthsafeguarding}. Recent large-scale systems like DeepSeek-V4~\cite{deepseekv4}, MiMo~\cite{xiao2026mimo}, and Qwen-3~\cite{qwen3technicalreport} demonstrate its efficacy as a post-training stage, while self-distillation studies~\cite{zhao2026self,shenfeld2026self,hubotter2026reinforcement} suggest OPD is a general mechanism for amplifying latent model capabilities rather than mere imitation. However, the success of OPD is often hindered by teacher-student misalignment. While prior research examines this incompatibility through global factors like vocabulary, reasoning style, or model family~\cite{fu2026revisiting,li2026rethinking}, our work focuses on a fine-grained source of stagnation. We investigate \textbf{Rock Tokens}—persistent high-loss positions that resist alignment even after training plateaus—to determine when OPD supervision provides productive signals versus unproductive noise.



\subsection{Token-Level Learning Dynamics in Reasoning Models}

Recent studies demonstrate that reasoning in LLMs is disproportionately supported by critical ``anchor'' tokens, the perturbation of which during inference significantly degrades model performance~\cite{gupta2025llms}. In parallel, RLVR research reveals that training is non-uniform~\cite{yang2025token,cheng2026reasoning,wu2025mitigating}, concentrating optimization on high-entropy "forking tokens" that act as decisive branching points for reasoning gains~\cite{wangbeyond,li2025faedkvinfinitewindowfouriertransform}, and further analysis confirms the overall token distribution drift mainly happens on some sparse but crucial tokens~\cite{meng2026sparse}. Although OPD's dense supervision suggests high-loss tokens drive training, their specific contribution mechanism remains unclear—especially when they persist as \textbf{Rock Tokens}. It is unknown whether these recalcitrant positions are indispensable reasoning pillars or merely unproductive residues. We address this gap by systematically analyzing their gradient dynamics and training effects to determine their true role in the OPD process.